% !PRINTED 560
\weberpage{560}
%<<CONT
so ist die Determinante $D$ der quadratischen Form (1) (Bd.\,I, \S.\,59)
\weq{5}{D = \Sigma \pm a_{1,1}a_{2,2}\ldots a_{n,n} = A^2.}

Die $\alpha_{i,k}$ sind dann gleichfalls reelle Zahlen, die auf unendlich viele Arten bestimmt werden k\"onnen. Wir nehmen eine beliebige, aber weiterhin unver\"anderliche Transformation (2), (3) an.

Das Gleichungssystem (3) m\"oge, nach den $x$ aufgel\"ost, ergeben:
\weq{6}{x_i = \beta_{1,i}\xi_1 + \cdots + \beta_{n,i}\xi_n.}

Wir wollen nun den Variablen $x_i$ nur ganze rationale Zahlwerthe beilegen. Dann ist klar, dass die Function $f$ nur verschwinden kann, wenn alle $x_i$ verschwinden, und dass ihr Werth, wenn wir diesen Fall ausschliessen, nicht unter einen gewissen endlichen Grenzwerth $M$ heruntersinken kann, den wir das \gesperrt{Minimum der Form} $f$ nennen. Denn sinkt der Werth von $f$ unter einen Werth $\alpha^2$ herunter, so m\"ussen nach (2) alle $\xi_i$ dem absoluten Werthe nach unter $\alpha$ herunter sinken, und nach (6) w\"urden f\"ur ein hinl\"anglich kleines $\alpha$ s\"ammtliche $x_i$ unter 1 herabgedr\"uckt werden und m\"ussten dann, da es ganze Zahlen sind, alle gleich Null sein.

Es kommt vor Allem darauf an, nicht sowohl das Minimum genau zu bestimmen, als eine m\"oglichst kleine obere Grenze daf\"ur zu finden, \"uber die es nicht hinausgehen kann.

Wir bedienen uns jetzt mit \gesperrt{Minkowski} einer Ausdrucksweise, die einer Geometrie im Raume von $n$ Dimensionen entnommen ist, die ausserordentlich einfach zu dem gew\"unschten Resultat f\"uhrt. F\"ur die F\"alle $n=2$, $n=3$ sind die dabei gebrauchten Ausdr\"ucke vollst\"andig anschaulich und allereinfachster Art. Wir empfehlen dem Leser, sich die nun auszuf\"uhrenden Schritte an diesen beiden speciellen F\"allen anschaulich klar zu machen. Bei gr\"osseren Werthen von $n$ geht freilich die eigentliche Anschauung verloren. Man hat dann diese Ausdr\"ucke als kurze Bezeichnung gewisser analytischer Verh\"altnisse anzusehen, wobei noch die Analogie mit dem gew\"ohnlichen Raume die Auffassung erleichtert. Es hat nicht die geringste Schwierigkeit, alle diese Ausdr\"ucke durch rein analytische Formeln zu ersetzen; es w\"urde darunter aber sehr die K\"urze des Ausdrucks und die Leichtigkeit der Auffassung leiden.

Wir bezeichnen also jetzt die Variablen $\xi_1, \ldots, \xi_n$ als Coor-
%CONT>>
